\begin{Definition}
\label{def:banach-valued-analytic-function}
Let $G\subset\mathbb{C}$ be a domain and let $Y$ be a Banach space. Then a mapping
\begin{equation}\label{eq:banach-valued-complex-function-map}
\mathbb{C}\supset G\ni\lambda\mapsto B(\lambda)\in Y
\end{equation}
is said to be \textbf{analytic} at a point $\lambda_0\in\mathbb{C}$ if
\begin{equation}\label{eq:banach-valued-analytic-function-difference-quotient}
\lim_{\lambda\to\lambda_0}\dfrac{B(\lambda)-B(\lambda_0)}{\lambda-\lambda_0}
\end{equation}
exists.
\end{Definition}

As we implied, our main result is the following.

\begin{Theorem}
\label{thm:bounded-operator-analytic-resolvent}
Let $A\in\mathcal{L}(X)$. Then the resolvent operator $R_\lambda(A)$ (thought of as a function of $\lambda$) is analytic on the resolvent set $\rho(A)$.
\end{Theorem}

\begingroup
\postdisplaypenalty=10000
\renewcommand{\endProof}{\par\nobreak\hfill$\blacksquare$\par\vspace{0.1cm}}
\begin{Proof}
The existence of the limit of the difference quotient follows directly form the power series representation shown in Theorem~\ref{thm:local-resolvent-power-series}.
\end{Proof}
\endgroup

We now assert that the techniques and results developed for analytic functions in a standard complex variables course can be used with impunity on analytic functions with values in a Banach space. For a more thorough development of this idea; see e.g., \cite{DS}. As an example of an application of old techniques in this new setting we now prove the following.

\begin{Theorem}
\label{thm:bounded-operator-nonempty-spectrum}
The spectrum of a bounded operator on a nonzero Banach space has at least one element.
\end{Theorem}

\begingroup
\postdisplaypenalty=10000
\renewcommand{\endProof}{\par\nobreak\hfill$\blacksquare$\par\vspace{0.1cm}}
\begin{Proof}
Let $A\in\mathcal{L}(X)$, and suppose $\sigma(A)$ is empty; i.e., the resolvent set is the entire complex plane. By Theorem~\ref{thm:bounded-operator-analytic-resolvent}, the resolvent operator $R_\lambda(A)$ (thought of as a function of $\lambda$) is entire; i.e., analytic on the entire complex plane. We now note that $\lambda\mapsto R_\lambda(A)$ is bounded on all of $\mathbb{C}$. To see this, note that by \eqref{eq:resolvent-series-outside-norm-disk} we can get
\begin{equation}\label{eq:resolvent-large-parameter-bound}
\|R_\lambda(A)\|\leq\dfrac{1}{\|A\|}\quad\text{for }|\lambda|\geq2\|A\|.
\end{equation}
In addition, $\lambda\mapsto R_\lambda(A)$ must be bounded on any bounded disk since it is analytic. Thus, we can use Liouville\textquotesingle s theorem to deduce that $\lambda\mapsto R_\lambda(A)$ is a constant. This is a contradiction and completes the proof.
\end{Proof}
\endgroup

\begin{Remark}
\label{rem:unbounded-operator-empty-spectrum-example}
Theorems~\ref{thm:local-resolvent-power-series} and \ref{thm:bounded-operator-analytic-resolvent} can be extended (with similar proofs) to closed operators (cf. Problem~\ref{ex:closed-unbounded-resolvent-analyticity}). However, it is possible for an unbounded operator to have an empty spectrum. For example, let $X=L^2(0,1)$ and let
\begin{equation}\label{eq:empty-spectrum-derivative-domain}
\mathcal{D}(S):=\{y\in H^1(0,1)\mid y(0)=0\}
\end{equation}
and
\begin{equation}\label{eq:empty-spectrum-derivative-action}
Sy=i\dfrac{\mathrm{d}y}{\mathrm{d}x}.
\end{equation}
The reader should verify that for any $\lambda\in\mathbb{C}$, the operator $L_\lambda$ given by
\begin{equation}\label{eq:empty-spectrum-derivative-explicit-resolvent}
L_\lambda(y)(x):=-i\int_0^xe^{-i\lambda(x-s)}y(s)\,\mathrm{d}s
\end{equation}
with domain
\begin{equation}\label{eq:empty-spectrum-derivative-resolvent-domain}
\mathcal{D}(L_\lambda):=L^2(0,1)
\end{equation}
is indeed the resolvent operator $R_\lambda(S)$.
\end{Remark}
